\documentclass[12pt]{article}
\usepackage[utf8]{inputenc}
\usepackage[T1]{fontenc}
\usepackage{hyperref}
\title{Terebu Language Resources for K-12 Curriculum Development}
\author{Compiled by Claude Sonnet 5 (claude-sonnet-5) for Andr\'as Kornai}
\date{}
\begin{document}
\maketitle
\section{Introduction}
Terebu (ISO 639-3 trb; also Turubu, Terepu, Turupu) is an Oceanic Austronesian language of the Kairiru branch. It is spoken in and around Turubu village on the Taul coast of East Sepik Province, Papua New Guinea. Wikipedia cites about 130 speakers (1990). Joshua Project gives a population of 500 for the Terepu people.
\section{Primary Educational Resources}
No primers, textbooks or school programmes were found in the search results.
\section{Linguistic and Reference Materials}
\begin{itemize}
\item Wikipedia, Terebu language: \url{https://en.wikipedia.org/wiki/Terebu_language}
\item Numeral data (incomplete), Max Planck EVA: \url{https://lingweb.eva.mpg.de/channumerals/Terebu.htm}
\item ASJP word list entry: \url{https://asjp.clld.org/languages/TEREBU}
\item Endangered Languages Project entry: \url{https://endangeredlanguages.com/lang/1811?hl=fr}
\end{itemize}
\section{Religious and Community Resources}
Joshua Project says Bible translation has been started but not finished. No completed scripture was confirmed. \url{https://joshuaproject.net/people_groups/15442}
\section{Digital and Audio Resources}
No audio, video or corpus resources were found.
\section{Curriculum Development Context}
Only wordlist-level and encyclopedic material exists. A curriculum would need new documentation work with the community, and no existing materials can be adapted for it.
\end{document}